\section{Requerimientos Funcionales}
\begin{table}[h!]
\centering
\small
\caption{Listado Consolidado de Requerimientos Funcionales.}
\label{tab:func_req}
\begin{tabular}{@{} l p{9.5cm} l @{}}
\toprule
\rowcolor[HTML]{2D3E50} 
\textbf{\color[HTML]{FFFFFF} ID} & \textbf{\color[HTML]{FFFFFF} Descripción del Requerimiento} & \textbf{\color[HTML]{FFFFFF} Actor} \\ \midrule
\rowcolor[HTML]{F2F4F7} 
RF-01 & Registro de queja con generación de folio único automático. & Quejoso \\
RF-02 & Validación de quejas, consulta de antecedentes y canalización. & Recepcionista \\
\rowcolor[HTML]{F2F4F7} 
RF-03 & Análisis de expedientes, notas internas y acuerdos de admisión. & Analista \\
RF-04 & Gestión de plazos de respuesta y actos de investigación. & Subdefensor \\
\rowcolor[HTML]{F2F4F7} 
RF-05 & Consulta de tableros de control de gestión global y estadísticas. & Defensor \\
RF-06 & Acceder al panel de quejas con acuerdo de conclusión pendiente. & Defensor \\
\rowcolor[HTML]{F2F4F7} 
RF-07 & Realizar la firma digital de acuerdos de conclusión de forma individual o masiva. & Defensor \\
RF-08 & Cambiar automáticamente el estatus de la queja a "Finalizado" tras la firma. & Defensor \\
\rowcolor[HTML]{F2F4F7} 
RF-09 & Enviar expedientes concluidos al archivo histórico del sistema. & Defensor \\
RF-10 & Asignar roles y permisos específicos a los usuarios de la plataforma. & Defensor \\ \bottomrule

\end{tabular}
\end{table}